\documentclass[12pt]{article}

% --- Encoding / language (XeLaTeX-safe) ---
\usepackage{fontspec}
\usepackage[russian,english]{babel}
\defaultfontfeatures{Ligatures=TeX}
\setmainfont{Times New Roman}

% --- Math / theorem ---
\usepackage{amsmath}
\usepackage{amssymb}
\usepackage{amsthm}
\usepackage{stmaryrd}
\usepackage{breqn}
\usepackage{euscript}
\usepackage{mathrsfs}
\usepackage{enumerate}
\usepackage{empheq}

% --- Graphics / layout ---
\usepackage{graphicx}
\usepackage{wrapfig}
\usepackage{subcaption}
\usepackage{float}
\usepackage{xcolor}
\usepackage{geometry}

\geometry{
    left=2.5cm,
    right=2.5cm,
    top=2cm,
    bottom=2cm,
    bindingoffset=0cm
}

% --- Extra utilities ---
\usepackage{pifont}
\usepackage{lipsum}
\usepackage{todonotes}
\usepackage{csquotes}

% --- Bibliography ---
\usepackage[backend=biber,style=alphabetic]{biblatex}
\addbibresource{references.bib}

% --- Hyperref (must be loaded last) ---
\usepackage{hyperref}

% --- Custom imports ---
\input{notations}

% --- Safe image inclusion (shows empty figure if file missing) ---
\newcommand{\safeincludegraphics}[2][]{%
    \IfFileExists{#2}{%
        \includegraphics[#1]{#2}%
    }{%
        \textcolor{gray}{\rule{0.8\linewidth}{0.6\linewidth}}%
    }%
}

% --- Theorems ---
\newtheorem{proposition}{Proposition}
\newtheorem{theorem}{Theorem}
\theoremstyle{remark}
\newtheorem{remark}{Remark}

% --- Colors ---
\definecolor{darkred}{RGB}{139,0,0}

% --- Date helper ---
\newcommand{\todayYYMMDD}{\the\year-\the\month-\the\day}


\begin{document}

\hfill
% \todayYYMMDD
\begin{center}
% \section*{Custom course or document title}
\textit{Quantitive finance homework 7 exercise 2}\\
% \textit{Optional subtitle}\\
\textit{Author: Polozkov Dmitrii, Spirin Stepan}\\
% \textit{Supervisor: Name}\\
\textit{Date: \today}
\end{center}

% \tableofcontents
% \newpage 

\section*{Exercise 2}
First of all, suppose American options are discussed. 

% Note that at expiration time $T$, the payoff of the option can be decomposed into:
% \begin{align}
%     \label{eq:decomp}
%     X^{call}:=(S_T-K)^+ &= (S_T-K)^+\mathbb{I}\{\max_{u\le T} S_u < H\} + (S_T-K)^+\mathbb{I}\{\max_{u\le T} S_u \ge H\}\\
%     &=:X^{call\_u-o}+X^{call\_u-i}\nonumber.
% \end{align}

First we motivate our arbitrage portfolio construction by the following observation.

Suppose at arbitrary time $t_0<T$, the (vanilla) american put option is traded at price $P_0$.
Notice as $P_0$ is a r.v., it can be decomposed into two expectations (purely formally):
\begin{align*}
    P_0 &= P_0\mathbb{I}\{\max_{u\le t} S_u < H\} + P_0\mathbb{I}\{\max_{u\le t} S_u \ge H\}\\
    &=:P_0^{U\&O}+P_0^{U\&I}.
\end{align*}

Hence---intuitively---if $P_0>P_0^{U\&O}+P_0^{U\&I}$, the vanilla option must be overpriced. 

Denote $\Delta P := P_0-P_0^{U\&O}-P_0^{U\&I}>0$.
Aiming at making an arbitrage, we short 1 unit of the vanilla one\footnote{which essentially means issuing this option} and long 1 units of both barrier options. 
Net CF at this moment is $\Delta P>0$.

We do not exercise the barrier options unless the vanilla one is exercised. 
Suppose it gets exercised at time $t_1\le T$. 
Immediately, exercise the barrier options.

To calculate CF at $t_1$, consider a random time $\tau=\inf\{t\;|\; S_t\ge H\}.$
If $\tau>t_1$, then the up-and-in barrier option never came into existence, i.e. it's price is zero. 
Clearly, the price of the other barrier option is equal to the price of the vanilla one, as the barrier condition becomes a deterministic boolean at time $t_1$. 
This implies that the symmetrical exercise of our barrier options exactly balance out the vanilla one getting exercised, i.e. the CF at $t_1$ is zero.
The case of $\tau\le t_1$ is analoguous. 

The total CF of the portfolio is the sum of CF at times $t_0$ and $t_1$, and it's equal to $\Delta P>0$ in all states of the world.

\end{document}

